\section{Low temperature and the cycle hierarchy}\label{sec:lowtemp}

When the rates are exponentially small in a temperature, the exit rates of
the faces are exponentially small too, and the competition between dimension
and exit rate takes place on the scale of the Freidlin and Wentzell cycles of
the chain. Away from a finite set of critical heights the winners are cycles
of least size among those deep enough to survive. At a critical height a
finite spectral problem decides, with noncycle winners on open windows.

The cycles, their depths and the exit times of a domain are classical. They
go back to Freidlin and Wentzell~\cite{freidlinwentzell2012}, and we use them
in the form of Catoni~\cite{catoni1999}, of Olivieri and
Vares~\cite[Chapter~6]{olivierivares2005} and of Cirillo, Nardi and
Sohier~\cite{cirillonardisohier2015}. The depth of a face
(Theorem~\ref{thm:depth}) is a short consequence of Catoni's theorem on the
exit time of a domain. The asymptotics of the exit rates with their
prefactors (Proposition~\ref{prop:prefactor}) are classical in substance, and
go back to Wentzell~\cite{wentzell1972} and to Freidlin and
Wentzell~\cite[Chapter~6]{freidlinwentzell2012}. What is new here is the score
that weighs the dimension of a face against its exit rate, the finite
competition of envelopes at a critical depth, the open window on which a
noncycle wins, and the reading of the cycle tree from the winners.

\subsection{Arrhenius families}\label{sec:arrhenius}

Let $V_{ij}\ge0$ for $i\ne j$ be fixed costs and let
\[
 q_{ij}(\eta)=e^{-V_{ij}/\eta},\qquad \eta>0,
\]
so every off-diagonal rate is positive at every positive temperature. We
write $Q_\eta$ for the generator, and for a face $F$ we write $\lambda_F(\eta)$
for the exit rate $\lambda_F$ of $F$ under $Q_\eta$. In this section the
letter $h$ is a height, not the step size $T/N$. For a proper face $F$ define
its depth
\[
 D(F)=\lim_{\eta\downarrow0}\bigl(-\eta\log \lambda_F(\eta)\bigr),
\]
whose existence is part of Theorem~\ref{thm:depth}. Put $D(V)=\infty$. A fixed
height $h>0$ sets the time scale
\[
 \tau_\eta=e^{h/\eta},\qquad T=\tau_\eta L.
\]
Throughout the section we take $N\to\infty$ first at fixed $\eta$ and then
$\eta\downarrow0$, and every threshold in $N$ may depend on $\eta$. The start
$\mu_\eta$ may depend on $\eta$ but has full support and does not depend on
$N$. Since all rates are positive, every face is irreducible, and with a
full-support start every face satisfies (H) by Lemma~\ref{lem:admissible}.
So Theorem~\ref{thm:first}(iii) applies at every fixed $\eta$ to every face.
At fixed $\eta$ the face maxima of Theorem~\ref{thm:facemax} give
\begin{equation}\label{eq:lowtemp-score}
 \log M_F=-\Phi_{\eta,h}(F)\,L+\frac{|F|-1}2\log L+O_\eta(1),\qquad
 \Phi_{\eta,h}(F)=|F|-1+e^{h/\eta}\lambda_F(\eta).
\end{equation}
So $\Phi_{\eta,h}$ is the exponent $\Phi_\tau$ of Theorem~\ref{thm:first} for the
generator $Q_\eta$ at $\tau=e^{h/\eta}$. The constants in~\eqref{eq:lowtemp-score}
may depend badly on $\eta$, and nothing below exchanges the two limits.

We use the cycles of Freidlin and Wentzell in Catoni's
convention~\cite[Definition~4.2]{catoni1999}. A nonempty
proper set $C$ is a cycle if it is a singleton or if, for all distinct
$x,y\in C$, the probability of leaving $C$ before hitting $y$, starting from
$x$, is exponentially small in $1/\eta$. Its depth $\Gamma(C)$ is the
exponential scale of its mean exit time, which does not depend on the
starting point in $C$~\cite[Proposition~4.6]{catoni1999}, and
$\Gamma(\{i\})=\min_{j\ne i}V_{ij}$.
For a reversible family write $V_{ij}=H_{ij}-U_i$ with $H$ symmetric, so that
$\pi_i(\eta)\propto e^{-U_i/\eta}$. Then a proper set $C$ with at least 2 states
is a cycle exactly when every two of its points communicate strictly below its
boundary height $b(C)=\min_{i\in C,j\notin C}H_{ij}$, that is, are joined by a
path whose saddles $H$ are all less than $b(C)$, and $\Gamma(C)=b(C)-\min_{i\in C}U_i$
\cite[Definition~3.5, Propositions~3.10 and~3.15]{cirillonardisohier2015}.
We call this the energy criterion. The inequality is strict, since a weak one
would admit sets at a common saddle level that are not cycles.

\begin{theorem}[the depth of a face]\label{thm:depth}
For every nonempty proper face $F$ the depth $D(F)$ exists and
\[
 D(F)=\max\{\Gamma(C):C\subseteq F\ \text{a cycle}\}.
\]
In particular $D(C)=\Gamma(C)$ for every cycle $C$, and $F$ contains a cycle
$C$ with $D(C)=D(F)$. Moreover, with $\mathcal G_F=(-Q_F)^{-1}$ and
$M^{\mathrm{exit}}_F=\max_{i\in F}\E_i\tau_{F^c}$,
\[
 \frac1{M^{\mathrm{exit}}_F}\le \lambda_F(\eta)\le\frac{|F|}{M^{\mathrm{exit}}_F},
\]
so that $D(F)=\lim_{\eta\downarrow0}\eta\log M^{\mathrm{exit}}_F$.
\end{theorem}

\begin{proof}
Fix $\eta$ and a proper face $F$, and write $\lambda_F$ for $\lambda_F(\eta)$.
By Lemma~\ref{lem:face-monotone} with $v=0$, $\lambda_F>\lambda_V=0$. So
$\mathcal G_F=\int_0^\infty e^{tQ_F}\,dt$ converges, and since
$(e^{tQ_F})_{ij}=\Prob_i(X_t=j,\ t<\tau_{F^c})$ for the chain $X$ with
generator $Q_\eta$,
\[
 (\mathcal G_F)_{ij}=\E_i\int_0^{\tau_{F^c}}\mathbf 1\{X_t=j\}\,dt
 =\Prob_i(\tau_j<\tau_{F^c})\,(\mathcal G_F)_{jj}\le(\mathcal G_F)_{jj}
\]
by the strong Markov property at $\tau_j$.

\emph{The exit-time bounds.} The row sums of $\mathcal G_F$ are the mean exit
times $\E_i\tau_{F^c}$. The spectral radius $\rho$ of a nonnegative matrix is at
most its largest row sum and at least each of its diagonal entries. Hence
\[
 \rho(\mathcal G_F)\le M^{\mathrm{exit}}_F\le\operatorname{tr}\mathcal G_F
 \le|F|\max_j(\mathcal G_F)_{jj}\le|F|\rho(\mathcal G_F).
\]
The Perron vector of $-Q_F$ is positive and is an eigenvector of
$\mathcal G_F$ for $1/\lambda_F$. A positive eigenvector of a nonnegative
matrix belongs to its spectral radius~\cite{bermanplemmons1994}, so
$\rho(\mathcal G_F)=1/\lambda_F$, which gives the bounds.

\emph{The depth.} Fix $c>d$ and put $P_\eta=I+Q_\eta/c$. Every exit rate is
less than $d$, so $P_\eta$ is stochastic with holding probabilities at least
$1-d/c>0$. Then $-\eta\log P_\eta(i,j)\to V_{ij}$ for $i\ne j$ and
$-\eta\log P_\eta(i,i)\to0$. So $(P_\eta)$ has rare transitions in the sense
of~\cite[Definition~2.1]{catoni1999}, with $\beta=1/\eta$ and the rate function
equal to $V_{ij}$ off the diagonal and to $0$ on it, and this rate function is
irreducible since all rates are positive. The chain $P_\eta$
visits the same states in the same order as $X$, so the two chains have the
same cycles. The continuous chain is the discrete one run at the jump times of
a Poisson clock of rate $c$, so its mean exit times are the discrete ones
divided by $c$. For a domain $D$ and $x\in D$, Catoni's
Proposition~4.2~\cite{catoni1999} gives the limit $H_D(x)$ of $\beta^{-1}\log$ of
the mean exit time from $D$ started at $x$, and his Proposition~4.11 states that
for every domain $D$ other than the whole space,
$\max_{x\in D}H_D(x)=\max\{\Gamma(C):C\in\mathcal M(D)\}$, where $\mathcal M(D)$
is the set of maximal cycles contained in $D$ (his Definition~4.3) and
$\Gamma(C)$ is the depth of his Proposition~4.6. With $D=F$ this reads
\[
 \lim_{\eta\downarrow0}\eta\log M^{\mathrm{exit}}_F=\max_{C\in\mathcal M(F)}\Gamma(C).
\] These
are cycles of the whole chain, not of a chain restricted to $F$. By the
exit-time bounds $-\eta\log\lambda_F$ differs from $\eta\log M^{\mathrm{exit}}_F$
by at most $\eta\log|F|$, so $D(F)$ exists and equals this limit. Every cycle
$C\subseteq F$ lies in some $C'\in\mathcal M(F)$, and $\Gamma(C)\le\Gamma(C')$
since $\tau_{C^c}\le\tau_{C'^c}$ pathwise. So the maximum over $\mathcal M(F)$
is the maximum over all cycles in $F$, attained at some $C\in\mathcal M(F)$. If
$F$ is a cycle, then $\mathcal M(F)=\{F\}$ and $D(F)=\Gamma(F)$.
\end{proof}

\begin{theorem}[noncritical heights]\label{thm:noncritical}
Let $\mathcal D$ be the finite set of depths of proper cycles and let
$h\notin\mathcal D$. Let $k(h)$ be the least size of a cycle of depth greater
than $h$, where $V$ counts as a cycle of infinite depth and size $d$. Then for
all sufficiently small $\eta$, every minimizer of $\Phi_{\eta,h}$ is a cycle of
size $k(h)$ and depth greater than $h$, and every eventual support of the
global mode, in the limit $N\to\infty$ at that $\eta$, is such a cycle.
\end{theorem}

\begin{proof}
By Theorem~\ref{thm:depth} the depth of every proper face lies in $\mathcal D$
and differs from $h$. If $D(F)<h$, then
$e^{h/\eta}\lambda_F(\eta)=e^{(h-D(F)+o(1))/\eta}\to\infty$. If $D(F)>h$, then
$e^{h/\eta}\lambda_F(\eta)\to0$. The face $V$ has the score $d-1$. Call a face
deep if it is $V$ or a proper face of depth greater than $h$.

A deep face of least size is a cycle. Indeed, a deep proper face $F$ contains a
cycle $C$ with $D(C)=D(F)>h$ by Theorem~\ref{thm:depth}, and $C$ is deep, so
$C=F$ if $F$ has least size. Every cycle of depth greater than $h$ is deep. So
the deep faces of least size are the cycles of size $k(h)$ and depth greater
than $h$.

There are finitely many faces, so for all sufficiently small $\eta$ every face
that is not deep has $\Phi_{\eta,h}>d$, and every deep face of size $k(h)$ has
$\Phi_{\eta,h}<k(h)$. A deep face of larger size has $\Phi_{\eta,h}\ge k(h)$,
since $\lambda_F\ge0$. Hence every minimizer is a cycle of size $k(h)$ and
depth greater than $h$, and Theorem~\ref{thm:first}(iii) at $\tau=e^{h/\eta}$
puts the support of every global mode among the minimizers for large $N$.
\end{proof}

Theorem~\ref{thm:noncritical} is a theorem about least cardinality. It does
not say that every cycle of size $k(h)$ and depth greater than $h$ wins, and
the cycles that are minimal for inclusion among the deep ones can have other
sizes and are then never selected. Among faces of equal size, minimizing
$\Phi_{\eta,h}$ is minimizing $\lambda_F(\eta)$ at fixed $\eta$, whatever $h$
is. So the deeper cycle wins first, and at equal depth the smaller prefactor of
Proposition~\ref{prop:prefactor}. An exponentially small difference of rates
still decides, because $N\to\infty$ first. If two scores agree exactly at some
$\eta$, the terms of~\eqref{eq:lowtemp-score} below $L$ decide. The
full-support start is needed, since otherwise some cycles are not admissible
faces.

\begin{proposition}[prefactors]\label{prop:prefactor}
For every proper face $F$ there is $a_F\in(0,\infty)$ with
$\lambda_F(\eta)=a_Fe^{-D(F)/\eta}(1+o(1))$, namely $a_F=1/\rho(B_F)$ with
$B_F=\lim_{\eta\downarrow0}e^{-D(F)/\eta}\mathcal G_F(\eta)$, and
\[
 D(F)=\nu\bigl(\det(-Q_F)\bigr)-\min_{i\in F}\nu\bigl(\det(-Q_{F\setminus\{i\}})\bigr),
\]
where $\nu(f)$ is the exponent $b$ of the first nonzero term $c\,e^{-b/\eta}$
of a finite sum of exponentials $f$, and the determinant of the empty matrix
is $1$.
\end{proposition}

\begin{proof}
Every minor of $-Q_F$ is a finite real combination of exponentials
$e^{-b/\eta}$. After equal exponents are collected, a combination that is not
identically zero has a first nonzero term $c\,e^{-\nu/\eta}$ and equals
$c\,e^{-\nu/\eta}(1+o(1))$. The determinant of $-Q_F$ is positive, since its
real eigenvalues are at least $\lambda_F>0$ and the others come in conjugate
pairs. The entries of $\mathcal G_F$ are positive by the identity in the proof
of Theorem~\ref{thm:depth}. By Cramer's rule each is a cofactor divided by
$\det(-Q_F)$, so
\[
 (\mathcal G_F)_{ij}=\beta_{ij}e^{e_{ij}/\eta}(1+o(1)),\qquad \beta_{ij}>0 .
\]
The diagonal cofactor of $j$ is $\det(-Q_{F\setminus\{j\}})$, so
$e_{jj}=\nu(\det(-Q_F))-\nu(\det(-Q_{F\setminus\{j\}}))$. The exit-time
bounds give $\max_j(\mathcal G_F)_{jj}\le1/\lambda_F\le|F|\max_j(\mathcal G_F)_{jj}$,
so $D(F)=\max_je_{jj}$, which is the determinant formula. Since
$(\mathcal G_F)_{ij}\le(\mathcal G_F)_{jj}$, we have $e_{ij}\le D(F)$ for all
$i,j$. Hence $e^{-D(F)/\eta}\mathcal G_F$ converges entrywise to the matrix
$B_F$ with entries $\beta_{ij}$ where $e_{ij}=D(F)$ and $0$ elsewhere. It has a
positive diagonal entry, so $\rho(B_F)>0$. The spectral radius is continuous
in the entries, and $\rho(e^{-D(F)/\eta}\mathcal G_F)=e^{-D(F)/\eta}/\lambda_F$.
Hence $e^{D(F)/\eta}\lambda_F\to1/\rho(B_F)$.
\end{proof}

The asymptotics of principal eigenvalues of matrices with entries of order
$e^{-V_{ij}/\eta}$, and their expressions through graphs and determinants, are due to
Wentzell~\cite{wentzell1972} and Freidlin and
Wentzell~\cite[Chapter~6]{freidlinwentzell2012}, and the proof above is a short
version for killed generators. The proposition uses the exact form of the rates. A family
that satisfies only the logarithmic limits $-\eta\log q_{ij}\to V_{ij}$ need not have finite
prefactors.

\begin{theorem}[critical windows]\label{thm:critical-window}
Fix a depth $\gamma\in\mathcal D$ and $x>0$, and let $T=xe^{\gamma/\eta}L$.
Define the limiting score
\[
 J_{\gamma,x}(F)=\begin{cases}+\infty,&D(F)<\gamma,\\ |F|-1+xa_F,&D(F)=\gamma,\\ |F|-1,&D(F)>\gamma,\end{cases}
\]
with $V$ in the last case. If $J_{\gamma,x}$ has a unique minimizer $F_x$, then
for all sufficiently small $\eta$, $F_x$ is the unique minimizer of the score
$|F|-1+xe^{\gamma/\eta}\lambda_F(\eta)$ and the eventual support of the global mode
as $N\to\infty$ at that $\eta$. The same holds, with one threshold for $\eta$,
uniformly for $x$ in a compact interval on which one face has a positive gap.
The lower envelope
$x\mapsto\min_FJ_{\gamma,x}(F)$ is concave and piecewise linear, and its
minimizers need not be cycles.
\end{theorem}

\begin{proof}
Let $S_\eta(F)=|F|-1+xe^{\gamma/\eta}\lambda_F(\eta)$, which is
$\Phi_\tau(F)$ for $Q_\eta$ at $\tau=xe^{\gamma/\eta}$. By
Proposition~\ref{prop:prefactor},
$e^{\gamma/\eta}\lambda_F(\eta)=a_Fe^{(\gamma-D(F))/\eta}(1+o(1))$, which tends to
$\infty$, $a_F$ or $0$ as $D(F)<\gamma$, $D(F)=\gamma$ or $D(F)>\gamma$, and
$S_\eta(V)=d-1$. On a compact interval $I\subset(0,\infty)$ the convergence
$S_\eta(F)\to J_{\gamma,x}(F)$ is uniform in $x$, divergence included, since
$x$ is bounded below on $I$.

Suppose that one face $F$ satisfies $J_{\gamma,x}(G)\ge J_{\gamma,x}(F)+g$ for
all $G\ne F$ and $x\in I$, with $g>0$. There are finitely many faces, so there
is $\eta_0$ with $S_\eta(G)>S_\eta(F)$ for all $G\ne F$, $x\in I$ and
$\eta<\eta_0$. A unique minimizer $F_x$ at a single $x$ has such a gap on a
small interval around $x$, because the scores are continuous in $x$.
Theorem~\ref{thm:first}(iii) at $\tau=xe^{\gamma/\eta}$ then gives the support
of the global mode, with a threshold in $N$ that may depend on $\eta$ and $x$.

The function $x\mapsto J_{\gamma,x}(F)$ is affine for every face with
$D(F)\ge\gamma$, and $J_{\gamma,x}(V)=d-1$ is finite. So the envelope is the
minimum of finitely many affine functions, which is concave and piecewise
linear. Example~\ref{ex:four-state} shows a minimizer that is not a cycle.
\end{proof}

At a tie of the limiting lines, Theorem~\ref{thm:critical-window} identifies
the candidates only, and the exact scores at each $\eta$ decide. Along the
envelope the slope $a_F$ of the active line decreases with $x$ and its
intercept $|F|-1$ increases, as for the hull of Theorem~\ref{thm:hull}.
Putting $h=\gamma$ and treating the faces of depth $\gamma$ as faces of greater
depth would discard the term $xa_F$, which is of order $1$ in the coefficient
of $L$.

\begin{example}[a noncycle winner on an open window]\label{ex:four-state}
Take $U=(0,1,1,1)$ and the symmetric saddles $H_{12}=2$,
$H_{13}=H_{23}=H_{34}=3$, $H_{14}=H_{24}=5$, so that with $t=e^{-1/\eta}$
\[
 (q_{ij})_{i\ne j}=\begin{pmatrix}*&t^2&t^3&t^5\\ t&*&t^2&t^4\\ t^2&t^2&*&t^2\\ t^4&t^4&t^2&*\end{pmatrix}.
\]
The only proper cycle with more than one state is $\{1,2\}$, of depth $3$. The
faces of depth $3$ are $\{1,2\}$, $\{1,2,3\}$ and $\{1,2,4\}$, with prefactors
$a_{12}=2$, $a_{123}=\frac23$ and $a_{124}=2$, so at $T=xt^{-3}L$ the limiting
scores are $1+2x$, $2+\frac23x$, $2+2x$ and $3$ for $V$. Hence the eventual
winner is $\{1,2\}$ for $0<x<\frac34$, the noncycle $\{1,2,3\}$ for
$\frac34<x<\frac32$, and $V$ for $x>\frac32$. At $h=3$ and $x=1$ the noncycle
beats every cycle. Its window has width $\log2$ in the rescaled height
$(h-3)/\eta$ and collapses onto the single height $h=3$, which is why it does
not contradict Theorem~\ref{thm:noncritical}.
\end{example}

\begin{proof}
Here $q_{ij}=t^{H_{ij}-U_i}$ and $\pi\propto(1,t,t,t)$. The exit rates are
\[
 q_1=t^2+t^3+t^5,\qquad q_2=t+t^2+t^4,\qquad q_3=3t^2,\qquad q_4=t^2+2t^4 .
\]

\emph{Cycles and depths.} The singletons have depths $\min_jV_{ij}=2,1,2,2$. We
test the sets with 2 or 3 states by the energy criterion. The set $\{1,2\}$
has boundary height $3$ and internal saddle $H_{12}=2<3$, so it is a cycle of
depth $3-U_1=3$. A set that contains exactly one of the states $1$ and $2$ has
the edge $H_{12}=2$ to its complement, and no two of its points communicate
below $2$. The remaining sets $\{3,4\}$, $\{1,2,3\}$ and $\{1,2,4\}$ have
boundary height $3$. In the first two the state $3$ is joined to the others
only by saddles of height $3$, and in the third the state $4$ is joined to the
others only by saddles of height $5$. So none is a cycle. By
Theorem~\ref{thm:depth}, a face that contains $\{1,2\}$ has depth $3$ and every
other proper face has depth at most $2$. So $\mathcal D=\{1,2,3\}$, and the
faces of depth $3$ are the three named.

\emph{Prefactors.} We apply Proposition~\ref{prop:prefactor}. For $\{1,2\}$,
$\det(-Q_{12})=q_1q_2-t^3=2t^4+O(t^5)$, and the first column of the adjugate is
$(q_2,t)=(t,t)+O(t^2)$, while the second column is $(t^2,q_1)=O(t^2)$. Hence
$t^3\mathcal G_{12}\to B_{12}$ with first column $(\frac12,\frac12)$ and second
column $0$, so $\rho(B_{12})=\frac12$ and $a_{12}=2$. For $\{1,2,3\}$,
\[
 -Q_{123}=\begin{pmatrix}t^2+t^3+t^5&-t^2&-t^3\\-t&t+t^2+t^4&-t^2\\-t^2&-t^2&3t^2\end{pmatrix}.
\]
Its determinant is $2t^6+O(t^7)$. The first column of its adjugate has the
leading terms $3t^3,3t^3,2t^3$, and every other entry of the adjugate is
$O(t^4)$. So $B_{123}$ has first column $(\frac32,\frac32,1)$ and zeros
elsewhere. Its eigenvalues are $\frac32$ and $0$, so $a_{123}=\frac23$. For
$\{1,2,4\}$ the determinant is $2t^6+O(t^7)$, the first column of the adjugate
is $(t^3,t^3,0)+O(t^4)$ and the other entries are $O(t^4)$. So $B_{124}$ has
first column $(\frac12,\frac12,0)$ and $a_{124}=2$. These leading terms are the
output of an exact computation in rational arithmetic.

\emph{The winners.} At $\gamma=3$ the faces of depth less than $3$ have
$J_{3,x}=\infty$. The lines are $1+2x$, $2+\frac23x$, $2+2x$ and $3$. The face
$\{1,2,4\}$ is always $1$ above $\{1,2\}$. Then $1+2x<2+\frac23x$ exactly when
$x<\frac34$, $2+\frac23x<3$ exactly when $x<\frac32$, and $1+2x<3$ exactly
when $x<1$. So the minimizer is unique and equals $\{1,2\}$ on $(0,\frac34)$,
$\{1,2,3\}$ on $(\frac34,\frac32)$ and $V$ on $(\frac32,\infty)$, and
Theorem~\ref{thm:critical-window} turns each into the eventual support of the
global mode, uniformly on compact subintervals. We claim nothing at the ties
$x=\frac34$ and $x=\frac32$.

\emph{The noncycle.} At $x=1$, that is at $h=3$, the score of $\{1,2,3\}$ is
$\frac83$. The cycles of depth less than $3$ have infinite limiting score, and
the only cycles of depth at least $3$ are $\{1,2\}$ and $V$, both of score $3$.
So the noncycle beats every cycle. In the rescaled height $s=(h-3)/\eta=\log x$ the window is
$\log\frac34<s<\log\frac32$, of width $\log2$. Its heights $h=3+\eta s$ tend to
$3$, so at every fixed $h\ne3$ the window is eventually left.
\end{proof}

The mechanism is that state $3$ costs one dimension but returns the chain to
$\{1,2\}$ before it is killed at state $4$, which lowers the leading exit rate
from $2t^3$ to $\frac23t^3$. On the middle interval this gain is worth more
than the dimension. The selected face is the support of the counts, not of the
limiting occupation fractions. For $\{1,2,3\}$ the right Perron vector tends to
$(1,1,\frac23)$ up to scale, and with $\pi\propto(1,t,t,t)$ the quasi-ergodic
law $\alpha^*_{123}$ tends to the point mass at state $1$.

\begin{example}[a branch switch]\label{ex:branch-switch}
Take $U=(2,0,0)$, $H_{12}=H_{13}=5$ and $H_{23}=1$. Then $\{2,3\}$ is a cycle
of depth $5$, the singleton depths are $3,1,1$, and the winners are $\{1\}$
for $0<h<3$, $\{2,3\}$ for $3<h<5$ and $V$ for $h>5$. The two winning cycles
are disjoint nodes of the same cycle tree, so the selector jumps between
branches and never selects the nodes $\{2\}$ and $\{3\}$.
\end{example}

\begin{proof}
With $t=e^{-1/\eta}$ the rates are $q_{12}=q_{13}=t^3$, $q_{21}=q_{31}=t^5$ and
$q_{23}=q_{32}=t$. So $\lambda_{\{1\}}=2t^3$, $\lambda_{\{2\}}=\lambda_{\{3\}}=t+t^5$
and $\lambda_{\{2,3\}}=t^5$. By the energy criterion $\{2,3\}$ is a cycle, and
$\{1,2\}$, $\{1,3\}$ are not, since each has an edge of height $1$ to its
complement. Hence $\mathcal D=\{1,3,5\}$. Theorem~\ref{thm:noncritical} gives
the winner $\{1\}$ on $(0,3)\setminus\{1\}$, since $2t^3<t+t^5$, the winner
$\{2,3\}$ on $(3,5)$, where $k(h)=2$, and $V$ above $5$. At $h=1$,
$\Phi_{\eta,1}(\{1\})=2t^2\to0$ while every other face has a score that tends
to at least $1$, so Theorem~\ref{thm:first}(iii) applies directly. We claim
nothing at $h=3$ and $h=5$.
\end{proof}

\subsection{Reading the cycle tree}\label{sec:cycle-tree}

We now ask the converse question of Section~\ref{sec:reconstruction} at low
temperature. Suppose that we observe only the winning face, as a function of
the height and of a field. How much of the cycle structure do we learn? The
unperturbed path shows the depth of every face that it selects and nothing
else. Suppressing the complement of a face does not show its depth either. A
calibrated reward on the face does. The depths then determine the cycles and
the cycle tree, and the critical windows determine the prefactors of the faces
on their envelopes.

A field is a vector $v=v(\eta)\in\R^d$, which may depend on $\eta$ but not on
$N$. As in Section~\ref{sec:observation} we write
$\lambda_F(v)=-\varrho(Q_{\eta,F}+\operatorname{diag}v_F)$, with $Q_{\eta,F}$ the
killed generator of $F$ under $Q_\eta$, and
\[
 \Phi_{\eta,h}(F;v)=|F|-1+e^{h/\eta}\lambda_F(v),
\]
so that $\Phi_{\eta,h}(F;0)=\Phi_{\eta,h}(F)$. In this subsection a height $h$
may be any real number. We say that a face $G$ \emph{wins at $(h,v)$} if there
is $\eta_0>0$ such that for every $\eta<\eta_0$ the face $G$ is the unique
minimizer of $\Phi_{\eta,h}(\cdot\,;v)$ over all nonempty faces. Where two
faces of equal size have identical rates, for instance two disjoint cycles
built from the same costs, no face wins and nothing is observed.

Proposition~\ref{prop:controlled-score} turns a win into an observation. At a
fixed $\eta$ it applies to $Q_\eta$, to the start $\mu_\eta$, to the fixed
vector $v(\eta)$ and to the single time scale $\tau=e^{h/\eta}$. Its proof uses
only results stated for an irreducible admissible face of a general generator,
and it allows a field of either sign. So if $G$ wins at $(h,v)$, then for every
$\eta<\eta_0$ the support of every maximizer of the weighted law is $G$ for all
large $N$. This is the only probabilistic input below.

Two elementary facts are used throughout. The Perron root of a Metzler matrix is
nondecreasing in its entries~\cite[Ch.~2]{bermanplemmons1994}, so $v\le w$
entrywise gives $\lambda_F(w)\le\lambda_F(v)$, and
$\lambda_F(v+\kappa\ones)=\lambda_F(v)-\kappa$ for every real $\kappa$. So
adding a constant to the field changes no winner. And by
Lemma~\ref{lem:face-monotone}, $\lambda_G(v)<\lambda_{F}(v)$ whenever
$F\subsetneq G$. In particular $D(F)\le D(G)$ for $F\subseteq G$.

\begin{proposition}[the unperturbed path]\label{prop:unperturbed}
Let $v=0$, and let a proper face $W$ win at some height $h_0\notin\mathcal D$.
Then $W$ is a cycle with $D(W)>h_0$, $W$ wins at every
$h\in[h_0,D(W))\setminus\mathcal D$, and $W$ wins at no $h>D(W)$. Hence the
supremum of the heights at which $W$ wins is $D(W)$.
\end{proposition}

\begin{proof}
By Theorem~\ref{thm:noncritical}, $W$ is a cycle of size $k(h_0)$ and depth
greater than $h_0$. Let $h\in[h_0,D(W))\setminus\mathcal D$. A face of depth
greater than $h$ has depth greater than $h_0$, so $k(h)\ge k(h_0)=|W|$, and
$k(h)=|W|$ since $D(W)>h$. By Theorem~\ref{thm:noncritical} again, for small
$\eta$ every minimizer of $\Phi_{\eta,h}$ is a cycle $G$ with $|G|=|W|$ and
$D(G)>h_0$. For such $G\ne W$,
$\Phi_{\eta,h}(G)-\Phi_{\eta,h}(W)=e^{h/\eta}(\lambda_G-\lambda_W)$ has the sign
of $\Phi_{\eta,h_0}(G)-\Phi_{\eta,h_0}(W)$, which is positive for small $\eta$.
Hence $W$ is the unique minimizer. For $h>D(W)$ we have
$e^{h/\eta}\lambda_W\to\infty$, while $\Phi_{\eta,h}(V)=d-1$.
\end{proof}

So the path shows the depth of every face it selects, as the supremum of the
heights at which it is selected. The noncritical heights at which a face wins
form an interval with finitely many critical heights removed, and these may be
depths of faces that never win, so only the supremum is $D(W)$. The path shows
nothing about the other faces, and a noncycle winner of a critical window
occupies a single height. In Example~\ref{ex:four-state} the singleton $\{4\}$
wins on $(0,1)\cup(1,2)$, since $q_4<q_1<q_3$ and $q_4<q_2$ for small $\eta$,
then $\{1,2\}$ wins on $(2,3)$ and $V$ above $3$. The cut at $1$ is the depth
of the cycle $\{2\}$, which never wins. The path shows $D(\{4\})=2$ and
$D(\{1,2\})=3$ and nothing about the cycles $\{1\}$, $\{2\}$ and $\{3\}$. It is
read against an absolute reference, since $\lambda_V(0)=0$ makes the score of
$V$ exactly $d-1$.

Killing the complement of a face at a large rate $M$ does not read its depth. The
killing gives $V$ an exit rate below $\lambda_F$, and the winner leaves $F$ only at the
height where the gap between $\lambda_F$ and that rate, the flux that returns to $F$
through the killed states, becomes visible. This gap can be of smaller order than
$\lambda_F$. For two states, $F=\{1\}$ and $M=e^{g/\eta}$, the exit rate $\lambda$ of $V$
is the smaller root of $(q_{12}-\lambda)(q_{21}+M-\lambda)=q_{12}q_{21}$, so the gap is
$q_{12}q_{21}/(q_{21}+M-\lambda)=e^{-(V_{12}+V_{21}+g)/\eta}(1+o(1))$, and the winner leaves
$\{1\}$ at the height $V_{12}+V_{21}+g$ and not at $D(\{1\})=V_{12}$.

The cure is to keep a face whose rate is known. For the unperturbed path this
is $V$, with $\lambda_V(0)=0$. Every field $v\ge0$ keeps a version of it, since
then $\lambda_V(v)\le0$ and $\Phi_{\eta,h}(V;v)\le d-1$. We therefore reward the
target face instead of killing its complement. For a proper face $G$, a fixed
constant $c>0$ and a real $\gamma$ put
\[
 v^{G,\gamma}=c\,e^{-\gamma/\eta}\ones_G .
\]
The field is constant on $G$, so $\lambda_G(v^{G,\gamma})=\lambda_G-c\,e^{-\gamma/\eta}$
exactly. Define the threshold of $G$ as
\[
 \vartheta(G)=\sup\bigl\{\gamma\in\R:\ G\text{ wins at }(h,v^{G,\gamma})\text{ for some }h>\gamma\bigr\},
\]
with $\sup\emptyset=-\infty$. It is read from the map $(h,v)\mapsto$ winner alone,
at heights where a winner exists. The observer needs no list of critical
heights, which are not known before the depths are.

The next two lemmas are standard facts of cycle theory. Laminarity is Catoni's
Proposition~4.4~\cite{catoni1999} and the tree of cycles of
Trouv\'e~\cite{trouve1996} and of Olivieri and Vares~\cite[Chapter~6]{olivierivares2005},
and the strict growth of depths along nested cycles can be read off Catoni's
Proposition~4.18. We include the short proofs for completeness, in the convention of
cycles through hitting probabilities used here.

\begin{lemma}[laminarity]\label{lem:laminar}
Two cycles are disjoint or nested.
\end{lemma}

\begin{proof}
Let $C_1,C_2$ be cycles with a common state $z$, and suppose that neither
contains the other. Take $x\in C_1\setminus C_2$ and $y\in C_2\setminus C_1$. Since
$y\notin C_1$, $\Prob_z(\tau_x<\tau_y)\ge\Prob_z(\tau_x<\tau_{C_1^c})$, which
tends to $1$ because $C_1$ is a cycle. Since $x\notin C_2$,
$\Prob_z(\tau_{C_2^c}<\tau_y)\ge\Prob_z(\tau_x<\tau_y)$, which tends to $1$ and
contradicts that $C_2$ is a cycle.
\end{proof}

\begin{lemma}[depth with a calibrated killing]\label{lem:killed-depth}
Let $K\subseteq[d]$, $\gamma\in\R$ and $s=c\,e^{-\gamma/\eta}$. For every proper
face $G$ the limit
$D^K(G)=\lim_{\eta\downarrow0}\bigl(-\eta\log\lambda_G(-s\ones_K)\bigr)$ exists, and
\[
 D^K(G)=\max\{D^K(C):C\subseteq G\ \text{a cycle}\}.
\]
Moreover $D^K(C)=D(C)$ if $C\cap K=\emptyset$, and $D^K(C)\le\gamma$ if
$C\subseteq K$.
\end{lemma}

\begin{proof}
Fix $\Lambda$ larger than every element of $\mathcal D$. Add a cemetery state
$\partial$. Every state of $K$ jumps to $\partial$ at rate $s$, every other state
of $[d]$ jumps to $\partial$ at rate $\varepsilon=e^{-\Lambda/\eta}$, and $\partial$
jumps to every state of $[d]$ at rate $\varepsilon$. All rates of this enlarged
chain are positive. If $\gamma<0$ we multiply all its rates by $e^{\gamma/\eta}$,
which shifts every depth by $-\gamma$ and changes no hitting probability, hence
no cycle. The proof of Theorem~\ref{thm:depth} uses the rates only through the
limits of $-\eta\log$ of the rates and through the exit-time bounds, once the
uniformization constant exceeds the largest total rate, here at most
$\max(d-1,c)+d\varepsilon$. So it applies to the enlarged chain.

For $G\subseteq[d]$ the killed generator of $G$ in the enlarged chain is
$Q_G-\operatorname{diag}(s\ones_{G\cap K}+\varepsilon\ones_{G\setminus K})$. By the two
elementary facts its exit rate $\lambda'_G$ satisfies
$\lambda_G(-s\ones_K)\le\lambda'_G\le\lambda_G(-s\ones_K)+\varepsilon$. Since
$\lambda_G(-s\ones_K)\ge\lambda_G$ and $D(G)<\Lambda$, we get
$\lambda'_G=\lambda_G(-s\ones_K)(1+o(1))$. So Theorem~\ref{thm:depth} for the
enlarged chain gives the existence of $D^K(G)$ and the identity
$D^K(G)=\max D^K(C)$ over the cycles $C\subseteq G$ of the enlarged chain.

Each such $C$ is a cycle of $Q_\eta$. Indeed, the enlarged chain is the chain
$X$ of $Q_\eta$ run until an independent killing time, and then $\partial$. If $X$
leaves $C$ before it hits $y$, so does the enlarged chain. So for distinct
$x,y\in C$ the probability $\Prob_x(\tau_{C^c}<\tau_y)$ for $X$ is at most the
same probability for the enlarged chain, which is exponentially small. Hence
$D^K(G)$ is at most the maximum of $D^K(C)$ over the cycles $C\subseteq G$ of
$Q_\eta$, and the reverse inequality holds by Lemma~\ref{lem:face-monotone}. If
$C\cap K=\emptyset$ then $\lambda_C(-s\ones_K)=\lambda_C$. If $C\subseteq K$ then
$\lambda_C(-s\ones_K)=\lambda_C+s\ge s$, so $D^K(C)\le\gamma$.
\end{proof}

\begin{lemma}[strict growth along nested cycles]\label{lem:nested}
If $C'\subsetneq C$ are cycles, then $D(C')<D(C)$. Hence $D(S)<D(C)$ for every
nonempty proper subset $S$ of a cycle $C$.
\end{lemma}

\begin{proof}
Fix $z\in C\setminus C'$. Since $C$ is a cycle there is $\alpha>0$ with
$\Prob_w(\tau_{C^c}<\tau_y)\le e^{-\alpha/\eta}$ for all distinct $w,y\in C$ and
all small $\eta$. For each $\eta$ choose $y\in C'$ with
$\E_y\tau_{C'^c}=M^{\mathrm{exit}}_{C'}$. Start at $y$ and cut the path into
rounds. A round runs until the chain leaves $C'$. If it then sits at a state
$w\in C\setminus C'$, the round continues until the chain returns to $y$ or
leaves $C$. A round ends outside $C$ only if the chain leaves $C'$ straight into
$C^c$, which forces $\tau_{C^c}<\tau_z$, or if it leaves $C$ from $w$ before it
reaches $y$. By the strong Markov property each round ends outside $C$ with
probability at most $2e^{-\alpha/\eta}$, and the rounds are independent and
identically distributed. So the number $R$ of rounds up to the exit from $C$ has
$\E R\ge e^{\alpha/\eta}/2$. Each round lasts at least its first part, of mean
$M^{\mathrm{exit}}_{C'}$, and $R$ is a stopping time for the rounds. Wald's
identity gives
\[
 M^{\mathrm{exit}}_C\ge\E_y\tau_{C^c}\ge\E R\cdot M^{\mathrm{exit}}_{C'}
 \ge\tfrac12e^{\alpha/\eta}M^{\mathrm{exit}}_{C'} ,
\]
and $D(C)\ge D(C')+\alpha$ by the last statement of Theorem~\ref{thm:depth}.
For the second claim, Theorem~\ref{thm:depth} gives a cycle $C''\subseteq S$
with $D(C'')=D(S)$, and $C''\subsetneq C$.
\end{proof}

\begin{theorem}[depths from a calibrated attraction]\label{thm:attraction}
Fix $c>0$.
\begin{enumerate}
\item[(a)] Let $G$ be a proper face. If $G$ wins at $(h,v^{G,\gamma})$ for some
$h>\gamma$, then $\gamma\le D(G)$. Hence $\vartheta(G)\le D(G)$.
\item[(b)] Let $C$ be a cycle, and put $D^-(C)=\max_{i\in C}D(C\setminus\{i\})$,
with $D^-(C)=-\infty$ for a singleton. Then $D^-(C)<D(C)$, and for every
$\gamma<D(C)$ the cycle $C$ wins at $(h,v^{C,\gamma})$ for every
$h\in(\max(\gamma,D^-(C)),D(C))$. Hence $\vartheta(C)=D(C)$.
\item[(c)] For every proper face $F$,
\[
 D(F)=\max_{G\subseteq F}\vartheta(G).
\]
\end{enumerate}
\end{theorem}

\begin{proof}
Write $s=c\,e^{-\gamma/\eta}$ and $\sigma=e^{h/\eta}$, and write $\Phi(G)$ for the
score of $G$ under the field at hand.

(a) Let $\gamma>D(G)$ and $h>\gamma$. Then $s$ is of smaller order than
$\lambda_G$, so $\sigma(\lambda_G-s)\ge\sigma\lambda_G/2\to\infty$. Since
$v^{G,\gamma}\ge0$, the score of $V$ is at most $d-1$, so $G$ is not a
minimizer for small $\eta$.

(b) The inequality $D^-(C)<D(C)$ is Lemma~\ref{lem:nested}. Fix
$h\in(\max(\gamma,D^-(C)),D(C))$ and write $v=v^{C,\gamma}$. The score of $C$ is
$|C|-1+\sigma\lambda_C-\sigma s$, where $\sigma\lambda_C\to0$ because $h<D(C)$ and
$\sigma s=c\,e^{(h-\gamma)/\eta}\to\infty$ because $h>\gamma$. On another face
$G$ the field is $s\ones_{G\cap C}=s\ones_G-s\ones_{G\setminus C}$, so
\[
 \lambda_G(v)=\lambda_G(-s\ones_{G\setminus C})-s .
\]
\emph{If $G\cap C=\emptyset$}, then $\lambda_G(v)=\lambda_G>0$, so $\Phi(G)\ge0$
while $\Phi(C)\to-\infty$.
\emph{If $G\subsetneq C$}, then $\lambda_G(v)-\lambda_C(v)=\lambda_G-\lambda_C$, and
$D(G)\le D^-(C)<h<D(C)$, so $\sigma(\lambda_G-\lambda_C)\ge\sigma\lambda_G/2\to\infty$.
\emph{If $G\supsetneq C$}, then
$\lambda_C(v)-\lambda_G(v)=\lambda_C-\lambda_G(-s\ones_{G\setminus C})\le\lambda_C$, so
$\Phi(G)-\Phi(C)\ge|G|-|C|-\sigma\lambda_C\to|G|-|C|\ge1$. The face $V$ is of this
kind.
\emph{Otherwise} $G$ meets $C$ without containing it or lying in it. Put
$K=G\setminus C$. Then $\lambda_G(v)-\lambda_C(v)=\lambda_G(-s\ones_K)-\lambda_C$. A
cycle $C'\subseteq G$ that met both $C$ and $K$ would be nested with $C$ by
Lemma~\ref{lem:laminar}, which is impossible since $C'$ meets $K$ and
$C\not\subseteq G$. So $C'$ lies in $G\cap C$ or in $K$. In the first case
Lemma~\ref{lem:killed-depth} gives $D^K(C')=D(C')\le D(G\cap C)\le D^-(C)$, and in
the second $D^K(C')\le\gamma$. By Lemma~\ref{lem:killed-depth},
$D^K(G)\le\max(D^-(C),\gamma)<h$, so $\sigma\lambda_G(-s\ones_K)\to\infty$, while
$\sigma\lambda_C\to0$.

In every case $\Phi(G)-\Phi(C)$ tends to $+\infty$ or has limit inferior at
least $1$. There are finitely many faces, so $C$ wins at $(h,v)$. The interval
of $h$ is nonempty for every $\gamma<D(C)$, so $\vartheta(C)\ge D(C)$, and (a) gives
equality.

(c) For $G\subseteq F$, (a) and the monotonicity of $D$ give
$\vartheta(G)\le D(G)\le D(F)$. By Theorem~\ref{thm:depth} there is a cycle
$C\subseteq F$ with $D(C)=D(F)$, and $\vartheta(C)=D(F)$ by (b).
\end{proof}

The theorem needs no assumption of uniqueness. Under its own attraction the
score of $C$ tends to $-\infty$ and every other score stays above it, so $C$
wins even when other faces tie, as for two disjoint cycles of equal rates. This
is the main advantage of the design over the unperturbed path. The theorem
uses only the threshold in $\gamma$ and not the right end of the heights at
which $C$ wins, since above $D(C)$ the faces that contain $C$ compete through
$\sigma(\lambda_C(v)-\lambda_G(v))$, which can diverge. A noncycle may win under its
own attraction, for instance in a window, and its threshold may reach $D(G)$.
This is why (c) takes the maximum over all subfaces and needs no prior
knowledge of which faces are cycles. In the deep-singleton family with
$V_{12}=V_{21}=1$, $V_{13}=V_{23}=3$ and $V_{31}=V_{32}=5$, the singleton
$\{3\}$ of depth $5$ wins on all of $(0,5)$ along the unperturbed path, so the
path never shows the cycle $\{1,2\}$ of depth $3$, while
Theorem~\ref{thm:attraction}(b) reads $\vartheta(\{1,2\})=3$.

\begin{theorem}[cycles from depths]\label{thm:cycles-from-depths}
For a nonempty proper face $C$ the following are equivalent.
\begin{enumerate}
\item[(i)] $C$ is a cycle.
\item[(ii)] $D(C')<D(C)$ for every nonempty proper subset $C'$ of $C$.
\item[(iii)] $|C|=1$, or $D(C\setminus\{i\})<D(C)$ for every $i\in C$.
\end{enumerate}
Moreover $\Gamma(C)=D(C)$ for every cycle. Hence the depth function determines
the set of cycles, their depths and the cycle tree, that is, the laminar
family of cycles ordered by inclusion with root $V$, and depths increase
strictly along every branch.
\end{theorem}

\begin{proof}
(i) implies (ii) by Lemma~\ref{lem:nested}, and (ii) and (iii) are equivalent
since $D$ is monotone. Suppose (ii). Theorem~\ref{thm:depth} gives a cycle
$C^*\subseteq C$ with $D(C^*)=D(C)$, and (ii) forces $C^*=C$. The identity
$\Gamma(C)=D(C)$ is in Theorem~\ref{thm:depth}, laminarity is
Lemma~\ref{lem:laminar}, and the strict growth is Lemma~\ref{lem:nested}.
\end{proof}

No nondegeneracy is assumed. Ties between the depths of different faces are
allowed, because the cycles are strict and Lemma~\ref{lem:nested} rules out a
tie along nested cycles. The weak form of the criterion fails. Take three
states with all $V_{ij}=1$. Then $\lambda_{\{1\}}=2e^{-1/\eta}$ and
$\lambda_{\{1,2\}}=e^{-1/\eta}$, so $D(\{1\})=D(\{2\})=D(\{1,2\})=1$. The set
$\{1,2\}$ is not a cycle, since from state $1$ the chain leaves it before it
hits $2$ with probability $\frac12$. Criterion (ii) rejects it, and the weak
criterion $D(C')\le D(C)$ accepts it. In Example~\ref{ex:four-state} the faces
$\{1,2,3\}$ and $\{1,2,4\}$ tie in depth with the cycle $\{1,2\}$ and are
rejected for the same reason. Together, Theorems~\ref{thm:attraction}(c)
and~\ref{thm:cycles-from-depths} say that the map $(h,v)\mapsto$ winner,
restricted to the fields $v^{G,\gamma}$, determines the depth of every face, the
cycles, their depths and the cycle tree.

The prefactors are read at a critical depth from the unperturbed path. Fix
$\gamma\in\mathcal D$, take $v=0$ and $T=xe^{\gamma/\eta}L$, and let
$J_{\gamma,x}$ be the limiting score of Theorem~\ref{thm:critical-window}. This
is the height $h=\gamma+\eta\log x$, which moves with $\eta$, so here we
observe the winner as a function of the window coordinate $x$. Write
$e(x)=\min_FJ_{\gamma,x}(F)$ for the lower envelope, and put $\bar a_G=a_G$ if
$D(G)=\gamma$ and $\bar a_G=0$ if $D(G)>\gamma$, so that
$J_{\gamma,x}(G)=|G|-1+x\bar a_G$ when $D(G)\ge\gamma$.

\begin{theorem}[prefactors from window endpoints]\label{thm:window-prefactors}
Suppose there are $0=x_0<x_1<\dots<x_m<x_{m+1}=\infty$ and faces $G_0,\dots,G_m$
with $G_{k-1}\ne G_k$ such that $G_k$ is the unique minimizer of $J_{\gamma,x}$
for every $x\in(x_k,x_{k+1})$.
\begin{enumerate}
\item[(a)] $D(G_m)>\gamma$, and $D(G_k)=\gamma$ for $k<m$.
\item[(b)] $\bar a_{G_m}=0$, and for $k=m,\dots,1$,
\[
 \bar a_{G_{k-1}}=\bar a_{G_k}+\frac{|G_k|-|G_{k-1}|}{x_k}.
\]
Hence the points $x_k$ and the sizes determine $a_{G_k}$ for every $k<m$.
\item[(c)] Let $G$ be a face of depth $\gamma$ other than $G_0,\dots,G_m$, and put
$s_G=\sup_{x>0}(e(x)-|G|+1)/x$. Then $a_G\ge s_G$, and every value
$a>s_G$ in place of $a_G$ gives the same unique minimizers. If $G'\subseteq G$
has depth $\gamma$, then $a_G\le a_{G'}$.
\end{enumerate}
\end{theorem}

\begin{proof}
By Theorem~\ref{thm:critical-window}, $G_k$ is the eventual winner, uniformly on
compact subintervals of $(x_k,x_{k+1})$, so the faces $G_k$ and the points $x_k$
are observed. The envelope is continuous, so the lines of $G_{k-1}$ and $G_k$
meet at $x_k$, that is,
$|G_{k-1}|-1+x_k\bar a_{G_{k-1}}=|G_k|-1+x_k\bar a_{G_k}$. This is the recursion
in (b). Since $G_{k-1}$ lies below $G_k$ to the left of $x_k$ and above it to
the right, $\bar a_{G_{k-1}}>\bar a_{G_k}\ge0$, and a face with $\bar a_G>0$ has
depth $\gamma$. As $x\to\infty$ every line of positive slope diverges, while
$J_{\gamma,x}(V)=d-1$. So the last piece has slope $0$, that is,
$D(G_m)>\gamma$ and $\bar a_{G_m}=0$. This proves (a) and (b). For (c), the
envelope is the minimum of the lines, so $|G|-1+xa_G\ge e(x)$ for all $x>0$,
which is $a_G\ge s_G$. If $a>s_G$, then $|G|-1+xa>e(x)$ for all $x>0$, and the unique
minimizers do not change. Finally $\lambda_G\le\lambda_{G'}$ for $G'\subseteq G$,
and at equal depths Proposition~\ref{prop:prefactor} turns this into
$a_G\le a_{G'}$.
\end{proof}

The reference is absolute because the last piece has slope exactly $0$, and
$V$, with $\lambda_V(0)=0$, guarantees a face with a flat line. In
Example~\ref{ex:four-state} at $\gamma=3$ the pieces are $G_0=\{1,2\}$,
$G_1=\{1,2,3\}$ and $G_2=V$, with $x_1=\frac34$, $x_2=\frac32$ and sizes
$2,3,4$. The recursion gives $a_{123}=(4-3)/\frac32=\frac23$ and
$a_{12}=\frac23+(3-2)/\frac34=2$, the values found in the example. The face
$\{1,2,4\}$ is never on the envelope. Here $s_{124}=\frac23$, attained on the
whole interval $(\frac34,\frac32)$, and $\{1,2\}\subseteq\{1,2,4\}$ gives
$a_{124}\le a_{12}=2$. So $\frac23\le a_{124}\le2$. At $a_{124}=\frac23$ the
line of $\{1,2,4\}$ would coincide with that of $\{1,2,3\}$ on the middle
interval, so the hypothesis of Theorem~\ref{thm:window-prefactors} excludes
the lower end, but the observed winners alone do not. The window thus recovers
$a_{12}=2$ and $a_{123}=\frac23$ and only bounds $a_{124}$, whose value is $2$.
A calibrated field does not obviously help, because a field changes the rate
of $V$ to a value that depends on $Q$, and the slope of $V$ is then unknown. In
this example the limit of $\pi$ and the limiting quasi-stationary laws of the
three faces of depth $3$ all sit at state $1$. So to first order a field
$c\,e^{-3/\eta}$ at state $1$ moves every line of the window, that of $V$
included, by $-cx$ and leaves the endpoints where they are, and a field of this
order at another state moves no line. Whether some other calibrated field
separates $a_{124}$ from $a_{12}$ is open.

\begin{remark}[what the design needs]\label{rem:design}
The attraction design needs three calibrated quantities. The first is the time
unit, so that the height $h$, equivalently $T=e^{h/\eta}L$, and the window
coordinate $x$ are known. The second is the amplitude $c\,e^{-\gamma/\eta}$ of
the field in the same unit. The third is the labelling of the states, which
places the field on $G$ and gives the sizes $|G|$. It needs no rate of $Q$, and
no step of this subsection uses reversibility. Theorem~\ref{thm:depth} rests on
Catoni's theorem and on the exit-time bounds, and the other lemmas use only the
Green matrix, the definition of a cycle and the strong Markov property.
Lemma~\ref{lem:killed-depth} applies Theorem~\ref{thm:depth} to a chain with a
cemetery, which is never reversible, so even for a reversible family the proof
of Theorem~\ref{thm:attraction} goes through the general theorem. The
selection step is Proposition~\ref{prop:controlled-score}, which holds for a
general generator.
\end{remark}

Section~\ref{sec:verify} records a check of Theorem~\ref{thm:attraction} with the
finite-$\eta$ scores on the $138$ proper faces of $11$ cost matrices,
reversible, non-reversible and tied, with no violation.
